\section{Cell Manifest, Versions and Sampling}
\label{app:manifest}

\begin{table}[h]
\centering\scriptsize
\caption{All cells, baseline prompt. Pilot cells (fewer than 30 runs) are
excluded from the reporting basis.}
\begin{tabular}{@{}lllrrrl@{}}
\toprule
Model & Tier & Domain & Runs & Att. & Fail & Status \\
\midrule
\code{claude-haiku-4-5} & medium & financebench & 32 & 263 & 186 &  \\
\code{claude-haiku-4-5} & medium & hotpotqa & 194 & 1{,}066 & 692 &  \\
\code{claude-opus-4-8} & medium & financebench & 30 & 101 & 11 &  \\
\code{claude-opus-4-8} & medium & hotpotqa & 54 & 164 & 2 &  \\
\code{claude-sonnet-4-6} & medium & hotpotqa & 50 & 178 & 28 &  \\
\code{claude-sonnet-5} & medium & financebench & 30 & 114 & 24 &  \\
\code{claude-sonnet-5} & medium & hotpotqa & 52 & 199 & 43 &  \\
\code{gpt-4o} & medium & hotpotqa & 70 & 250 & 40 &  \\
\code{gpt-4o-mini} & medium & financebench & 5 & 28 & 13 & pilot (excluded) \\
\code{gpt-4o-mini} & medium & hotpotqa & 70 & 242 & 33 &  \\
\code{gpt-5.6-luna} & medium & financebench & 30 & 97 & 7 &  \\
\code{gpt-5.6-luna} & medium & hotpotqa & 55 & 169 & 4 &  \\
\code{gpt-5.6-sol} & medium & financebench & 30 & 96 & 6 &  \\
\code{gpt-5.6-sol} & medium & hotpotqa & 52 & 159 & 3 &  \\
\code{gpt-5.6-terra} & medium & financebench & 30 & 98 & 8 &  \\
\code{gpt-5.6-terra} & medium & hotpotqa & 54 & 168 & 6 &  \\
\code{claude-sonnet-4-6} & medium\_v2 & hotpotqa & 100 & 324 & 24 &  \\
\code{gpt-5.6-luna} & medium\_v3\_loose\_regex & hotpotqa & 2 & 6 & 0 & pilot (excluded) \\
\code{gpt-5.6-sol} & medium\_v3\_loose\_regex & hotpotqa & 50 & 155 & 5 &  \\
\code{gpt-5.6-terra} & medium\_v3\_loose\_regex & hotpotqa & 50 & 154 & 4 &  \\
\bottomrule
\end{tabular}

\end{table}

\paragraph{Sampling.}
HotpotQA questions are the first records of the validation split in the
distractor configuration, in dataset order; FinanceBench questions are the
first records of the train split, with evidence truncated to 2{,}000
characters and a truncation marker inserted. There is no shuffle and no
seed.

\paragraph{Resolved versions.}
Anthropic models and the GPT-5.6 family return the requested alias
unchanged. \code{gpt-4o} resolves to \code{gpt-4o-2024-08-06} and
\code{gpt-4o-mini} to \code{gpt-4o-mini-2024-07-18}. Version capture was
added after the first sweeps: \pn{versions.gpt-4o.precapture} of
\pn{versions.gpt-4o.attempts} \code{gpt-4o} attempts and
\pn{versions.gpt-4o-mini.precapture} of \pn{versions.gpt-4o-mini.attempts}
\code{gpt-4o-mini} attempts on the reporting basis predate it.

\paragraph{Repeated questions.}
The Haiku~4.5 HotpotQA baseline cell contains \pn{robust.cell.runs} runs
over \pn{robust.cell.questions} distinct questions. Restricting it to the
earliest run of each question gives a failure rate of
\pn{robust.one_run.rate} (\pn{robust.one_run.failures} of
\pn{robust.one_run.attempts}) against \pn{robust.cell.rate}
(\pn{robust.cell.failures} of \pn{robust.cell.attempts}) for the full cell.

\paragraph{Zero-length completion.}
One attempt (\pn{zero_length.model}, \pn{zero_length.node} node) returned
an empty completion with \pn{zero_length.output_tokens} output tokens,
equal to the configured ceiling. It is classified \cat{parse\_error}. The
stop reason was not logged; the record is consistent with the output
budget being exhausted before any text block was emitted.

\section{Schemas and Prompts}
\label{app:schemas}

The three schema tiers and three prompt variants are reproduced from the
released code.

\IfFileExists{code/medium.py}{\lstinputlisting[language=Python,caption={Medium tier schema}]{code/medium.py}}{\emph{[Copy src/schemaeval/schemas/medium.py to paper/code/medium.py.]}}

\IfFileExists{code/prompts.py}{\lstinputlisting[language=Python,caption={Prompt variants}]{code/prompts.py}}{\emph{[Copy src/schemaeval/prompts.py to paper/code/prompts.py.]}}

The medium\_v2 tier differs only in the three character caps and the
medium\_v3\_loose\_regex tier only in \code{CANONICAL\_ID\_PATTERN}.

\section{Classifier Validation}
\label{app:classifier}

\paragraph{Provenance.}
Hand labelling was performed against an earlier classifier in which
Pydantic error-type mapping preceded the behavioural predicates.
Schema-mirrored outputs were consequently labelled \cat{missing\_field}
and prose-preamble outputs \cat{parse\_error}. We corrected the order and
reclassified the full corpus from stored completions and error details;
the released \code{errors.py} and \code{reclassify\_all.py} reproduce
every label.

\paragraph{Agreement.}
The sample was drawn from all \pn{full.failures} failures in the corpus,
including pilot cells and prompt variants. \Cref{tab:validation} reports
precision and recall of the corrected classifier against the \emph{blind}
human labels. The four disagreements were adjudicated as labelling
errors: three outputs labelled \cat{enum\_violation} carried a
\code{string\_pattern\_mismatch} on an identifier, and one labelled
\cat{markdown\_wrap} contained no fence. We headline blind agreement
because adjudication uses evidence the labeller did not consult.

\begin{table}[h]
\centering\small
\caption{Corrected classifier against blind human labels
($n=\pn{validation.n}$), full corpus frame.}
\label{tab:validation}
\begin{tabular}{@{}lrll@{}}
\toprule
Category & Predicted & Precision & Recall \\
\midrule
\cat{markdown\_wrap} & 25 & 25/25 & 25/26 \\
\cat{length\_bound} & 30 & 30/30 & 30/30 \\
\cat{regex\_pattern} & 30 & 27/30 & 27/27 \\
\cat{schema\_mirroring} & 42 & 42/42 & 42/42 \\
\cat{prose\_preamble} & 9 & 8/9 & 8/8 \\
\cat{enum\_violation} & 8 & 8/8 & 8/11 \\
\bottomrule
\end{tabular}

\end{table}

\paragraph{Corpus screens.}
Because the sample is stratified by predicted label, it cannot establish
corpus recall. For the Stage-1 categories we screened every failure in
the full corpus with predicates broader than the detectors. A fence at
any position matches \pn{screen.fence_any.matches} failures, of which
\pn{screen.fence_any.markdown_wrap} are classified \cat{markdown\_wrap}
and \pn{screen.fence_any.prose_preamble} \cat{prose\_preamble} (prose
before a later fence). A \code{"properties"} key matches
\pn{screen.properties.matches} failures, of which
\pn{screen.properties.schema_mirroring} are classified
\cat{schema\_mirroring}; this screen tests one meta-keyword where the
detector requires two, so it bounds recall only among failures containing
that key. An opening character other than a brace, bracket or fence
matches \pn{screen.nonjson_opener.matches} non-empty failures, of which
\pn{screen.nonjson_opener.prose_preamble} are \cat{prose\_preamble}.

\section{Identifier Perturbation}
\label{app:perturbation}

\begin{table}[h]
\centering\small
\caption{Sub-causes of violating identifiers in \cat{regex\_pattern}
failures, gpt-5.6-sol and gpt-5.6-terra, HotpotQA. Counts are
identifiers, not attempts.}
\label{tab:perturbation}
\begin{tabular}{@{}lrr@{}}
\toprule
Sub-cause & Baseline pattern & Loosened pattern \\
\midrule
Leading digit & 4 & 0 \\
Length overflow & 4 & 10 \\
Unicode contamination & 0 & 1 \\
Other character & 0 & 0 \\
\midrule
Violating identifiers & 8 & 11 \\
\bottomrule
\end{tabular}

\end{table}

The baseline tier uses \idpattern{} and the perturbed tier \idpatternloose{};
both retain the 31-character ceiling, and the loosened pattern also
governs the supporting-identifier check inside the answer validator.
Sub-causes are assigned mechanically per identifier: non-ASCII characters,
a leading digit not admitted by the first-character class, other
characters outside the class, and length outside 3--31. The cells hold
\pn{pert.medium.gpt-5.6-sol.failures} and
\pn{pert.medium.gpt-5.6-terra.failures} failures under the baseline
pattern and \pn{pert.medium_v3_loose_regex.gpt-5.6-sol.failures} and
\pn{pert.medium_v3_loose_regex.gpt-5.6-terra.failures} under the loosened
pattern, too few to support rates or shares. The single Unicode case
under the loosened pattern is an otherwise lowercase band name with a
non-Latin combining mark inserted.

\section{Per-Domain and First-Attempt Breakdowns}
\label{app:breakdowns}

\begin{table}[h]
\centering\scriptsize
\caption{Per-cell results, reporting basis.}
\begin{tabular}{@{}lllrrrl@{}}
\toprule
Model & Tier & Domain & Att. & Rate & Att./run & Dominant \\
\midrule
\code{claude-haiku-4-5} & medium & financebench & 263 & 70.7\% & 8.22 & \cat{markdown\_wrap} \\
\code{claude-haiku-4-5} & medium & hotpotqa & 1{,}066 & 64.9\% & 7.96 & \cat{markdown\_wrap} \\
\code{claude-opus-4-8} & medium & financebench & 101 & 10.9\% & 3.37 & \cat{length\_bound} \\
\code{claude-opus-4-8} & medium & hotpotqa & 164 & 1.2\% & 3.04 & \cat{regex\_pattern} \\
\code{claude-sonnet-4-6} & medium & hotpotqa & 178 & 15.7\% & 3.56 & \cat{length\_bound} \\
\code{claude-sonnet-4-6} & medium\_v2 & hotpotqa & 324 & 7.4\% & 3.24 & \cat{regex\_pattern} \\
\code{claude-sonnet-5} & medium & financebench & 114 & 21.1\% & 3.80 & \cat{length\_bound} \\
\code{claude-sonnet-5} & medium & hotpotqa & 199 & 21.6\% & 3.83 & \cat{markdown\_wrap} \\
\code{gpt-4o} & medium & hotpotqa & 250 & 16.0\% & 3.57 & \cat{schema\_mirroring} \\
\code{gpt-4o-mini} & medium & hotpotqa & 242 & 13.6\% & 3.46 & \cat{regex\_pattern} \\
\code{gpt-5.6-luna} & medium & financebench & 97 & 7.2\% & 3.23 & \cat{regex\_pattern} \\
\code{gpt-5.6-luna} & medium & hotpotqa & 169 & 2.4\% & 3.07 & \cat{regex\_pattern} \\
\code{gpt-5.6-sol} & medium & financebench & 96 & 6.2\% & 3.20 & \cat{regex\_pattern} \\
\code{gpt-5.6-sol} & medium & hotpotqa & 159 & 1.9\% & 3.06 & \cat{regex\_pattern} \\
\code{gpt-5.6-sol} & medium\_v3\_loose\_regex & hotpotqa & 155 & 3.2\% & 3.10 & \cat{regex\_pattern} \\
\code{gpt-5.6-terra} & medium & financebench & 98 & 8.2\% & 3.27 & \cat{regex\_pattern} \\
\code{gpt-5.6-terra} & medium & hotpotqa & 168 & 3.6\% & 3.11 & \cat{regex\_pattern} \\
\code{gpt-5.6-terra} & medium\_v3\_loose\_regex & hotpotqa & 154 & 2.6\% & 3.08 & \cat{regex\_pattern} \\
\bottomrule
\end{tabular}

\end{table}

\begin{table}[h]
\centering\scriptsize
\caption{First-attempt validity by node. The denominator is the number of
runs that reached the node, which depends on upstream success; the
identify node is the only unconditioned comparison.}
\begin{tabular}{@{}lllrr@{}}
\toprule
Model & Tier & Node & Reached & First-attempt valid \\
\midrule
\code{claude-haiku-4-5} & medium & answer & 143 & 0.0\% \\
\code{claude-haiku-4-5} & medium & gather & 166 & 0.0\% \\
\code{claude-haiku-4-5} & medium & identify & 166 & 0.0\% \\
\code{claude-opus-4-8} & medium & answer & 84 & 89.3\% \\
\code{claude-opus-4-8} & medium & gather & 84 & 100.0\% \\
\code{claude-opus-4-8} & medium & identify & 84 & 96.4\% \\
\code{claude-sonnet-4-6} & medium & answer & 50 & 84.0\% \\
\code{claude-sonnet-4-6} & medium & gather & 50 & 82.0\% \\
\code{claude-sonnet-4-6} & medium & identify & 50 & 82.0\% \\
\code{claude-sonnet-4-6} & medium\_v2 & answer & 100 & 96.0\% \\
\code{claude-sonnet-4-6} & medium\_v2 & gather & 100 & 96.0\% \\
\code{claude-sonnet-4-6} & medium\_v2 & identify & 100 & 84.0\% \\
\code{claude-sonnet-5} & medium & answer & 82 & 86.6\% \\
\code{claude-sonnet-5} & medium & gather & 82 & 76.8\% \\
\code{claude-sonnet-5} & medium & identify & 82 & 64.6\% \\
\code{gpt-4o} & medium & answer & 70 & 58.6\% \\
\code{gpt-4o} & medium & gather & 70 & 100.0\% \\
\code{gpt-4o} & medium & identify & 70 & 84.3\% \\
\code{gpt-4o-mini} & medium & answer & 70 & 85.7\% \\
\code{gpt-4o-mini} & medium & gather & 70 & 100.0\% \\
\code{gpt-4o-mini} & medium & identify & 70 & 78.6\% \\
\code{gpt-5.6-luna} & medium & answer & 85 & 96.5\% \\
\code{gpt-5.6-luna} & medium & gather & 85 & 100.0\% \\
\code{gpt-5.6-luna} & medium & identify & 85 & 90.6\% \\
\code{gpt-5.6-sol} & medium & answer & 82 & 98.8\% \\
\code{gpt-5.6-sol} & medium & gather & 82 & 100.0\% \\
\code{gpt-5.6-sol} & medium & identify & 82 & 90.2\% \\
\code{gpt-5.6-sol} & medium\_v3\_loose\_regex & answer & 50 & 98.0\% \\
\code{gpt-5.6-sol} & medium\_v3\_loose\_regex & gather & 50 & 100.0\% \\
\code{gpt-5.6-sol} & medium\_v3\_loose\_regex & identify & 50 & 92.0\% \\
\code{gpt-5.6-terra} & medium & answer & 84 & 98.8\% \\
\code{gpt-5.6-terra} & medium & gather & 84 & 100.0\% \\
\code{gpt-5.6-terra} & medium & identify & 84 & 84.5\% \\
\code{gpt-5.6-terra} & medium\_v3\_loose\_regex & answer & 50 & 100.0\% \\
\code{gpt-5.6-terra} & medium\_v3\_loose\_regex & gather & 50 & 100.0\% \\
\code{gpt-5.6-terra} & medium\_v3\_loose\_regex & identify & 50 & 92.0\% \\
\bottomrule
\end{tabular}

\end{table}

\section{Model-Level Constraints}
\label{app:crossfield}

Model validators run only after every field-level constraint passes.
On the reporting basis, \pn{crossfield.reached} attempts reached them:
\pn{crossfield.validated} that validated and \pn{modellevel.violations}
that passed field validation and failed the per-item
supporting-identifier check in the answer validator. That check follows
the confidence-bucket check in the same method, so those attempts also
passed it. No attempt matched any of the four cross-field messages
(\pn{crossfield.violations} violations). The remaining
\pn{crossfield.not_reached} failures did not reach the validators; many
failed during JSON parsing. Each node tests only its own checks: two
cross-field checks in identify, one in gather, and one cross-field plus
one per-item pattern check in answer.

\section{Divergence Metric}
\label{app:divergence}

For each pair of model--tier groups and each category we take the
absolute difference in the share of the group's failures held by that
category. Groups pool task domains and use the reporting basis; of
\pn{divergence.groups_total} groups with any failure,
\pn{divergence.groups_qualifying} have at least 20 failures, giving
\pn{divergence.pairs} pairs. The two largest values arise on different
categories from opposite vendors, \pn{divergence.top.0.pp} and
\pn{divergence.runner_up_category_pp} points, and differ by
\pn{divergence.gap_pp} points, so we attach no significance to their order. Without the threshold the maximum
is \pn{divergence.unrestricted.top.0.pp} points, but it involves a group
with very few failures.

\begin{table}[h]
\centering\small
\caption{Largest category-share differences among qualifying groups.}
\begin{tabular}{@{}lllr@{}}
\toprule
Category & Group A & Group B & pp \\
\midrule
markdown\_wrap & \code{claude-haiku-4-5@medium} & \code{claude-sonnet-4-6@medium} & 72.9 \\
markdown\_wrap & \code{claude-haiku-4-5@medium} & \code{claude-sonnet-4-6@medium\_v2} & 72.9 \\
markdown\_wrap & \code{claude-haiku-4-5@medium} & \code{gpt-4o@medium} & 72.9 \\
markdown\_wrap & \code{claude-haiku-4-5@medium} & \code{gpt-4o-mini@medium} & 72.9 \\
schema\_mirroring & \code{claude-haiku-4-5@medium} & \code{gpt-4o@medium} & 72.5 \\
schema\_mirroring & \code{claude-sonnet-4-6@medium} & \code{gpt-4o@medium} & 72.5 \\
schema\_mirroring & \code{claude-sonnet-4-6@medium\_v2} & \code{gpt-4o@medium} & 72.5 \\
schema\_mirroring & \code{claude-sonnet-5@medium} & \code{gpt-4o@medium} & 72.5 \\
\bottomrule
\end{tabular}

\end{table}
